\documentclass[11pt]{article}
\usepackage[a4paper,margin=2.5cm]{geometry}
\usepackage{amsmath}
\usepackage{amssymb}
\usepackage{enumitem}
\usepackage{parskip}
\setlist[enumerate,2]{label=(\alph*), itemsep=1pt, topsep=2pt}
\setlength{\parindent}{0pt}
\newcommand{\correct}[1]{\textbf{#1}\;$\checkmark$}
\begin{document}
\section*{Routing games}
\begin{enumerate}[itemsep=6pt]
  \item A Nash flow is a flow in which ...
  \begin{enumerate}
    \item the total cost is minimised
    \item every path carries equal flow
    \item all traffic uses a single path
    \item no individual can reduce their own cost by switching to another path
  \end{enumerate}
  \item An optimal flow is the flow that ...
  \begin{enumerate}
    \item equals the Nash flow in every game
    \item uses the most paths
    \item no driver can improve on individually
    \item minimises the total cost across all traffic
  \end{enumerate}
  \item The price of anarchy measures ...
  \begin{enumerate}
    \item the number of paths in the network
    \item the cost of the cheapest path
    \item the ratio of the Nash flow cost to the optimal flow cost
    \item how many drivers there are
  \end{enumerate}
  \item Braess's paradox is the observation that ...
  \begin{enumerate}
    \item the Nash flow always equals the optimal flow
    \item removing a road always helps
    \item adding a new road can make every driver's commute worse
    \item drivers never reach an equilibrium
  \end{enumerate}
  \item A feasible flow that minimises the potential function \(\Phi\) is ...
  \begin{enumerate}
    \item the optimal flow
    \item a flow that uses no congested edges
    \item never an equilibrium
    \item a Nash flow
  \end{enumerate}
  \item The optimal flow of a routing game equals the Nash flow of the game with ...
  \begin{enumerate}
    \item all costs set to zero
    \item the costs doubled
    \item each cost function replaced by its marginal cost
    \item a single combined path
  \end{enumerate}
  \item A feasible flow is one in which ...
  \begin{enumerate}
    \item the total cost is minimised
    \item the flow on the paths of each commodity sums to its demand
    \item no path is used
    \item all flow uses a single path
  \end{enumerate}
  \item Edge cost (latency) functions in a routing game are assumed to be ...
  \begin{enumerate}
    \item negative
    \item constant
    \item non-negative, continuous and non-decreasing in the flow
    \item decreasing in the flow
  \end{enumerate}
  \item The price of anarchy is always ...
  \begin{enumerate}
    \item at least 1
    \item negative
    \item exactly 1
    \item at most 1
  \end{enumerate}
  \item In Pigou's example (one fixed-cost route and one congestible route) the price of anarchy is ...
  \begin{enumerate}
    \item \(1\)
    \item \(4/3\)
    \item \(8/7\)
    \item \(2\)
  \end{enumerate}
  \item Braess's paradox arises because each driver ...
  \begin{enumerate}
    \item cannot find the shortest path
    \item is forced onto a single route
    \item pays a toll
    \item ignores the extra congestion they impose on others (an externality)
  \end{enumerate}
  \item At a Nash flow, all paths a commodity actually uses ...
  \begin{enumerate}
    \item are the cheapest in the whole network
    \item carry equal flow
    \item have different costs
    \item have equal (and minimal) cost
  \end{enumerate}
  \item The total cost of a flow is ...
  \begin{enumerate}
    \item \(\sum_{e} c_e(f_e) f_e\), summed over the edges
    \item the cost of the cheapest edge
    \item the price of anarchy
    \item the number of paths used
  \end{enumerate}
  \item The optimal flow is found by ...
  \begin{enumerate}
    \item maximising the total cost
    \item letting each driver choose selfishly
    \item using only one path
    \item minimising the total cost over all feasible flows
  \end{enumerate}
  \item Marginal-cost pricing aligns the selfish and social outcomes because ...
  \begin{enumerate}
    \item it sets all costs to zero
    \item it forbids the use of congested edges
    \item it removes all congestion
    \item it makes each driver internalise the congestion they cause
  \end{enumerate}
  \item A routing game is written \(( G, r, c)\), where \(r\) denotes the ...
  \begin{enumerate}
    \item number of edges in the network
    \item price of anarchy
    \item reciprocal of the cost
    \item amount of traffic (demand) each source-sink commodity must route
  \end{enumerate}
  \item The potential function of a routing game is ...
  \begin{enumerate}
    \item \(\Phi(f) = \sum_{e} \int_0^{f_e} c_e(x)\,dx\)
    \item \(\Phi(f) = \sum_e f_e\)
    \item \(\Phi(f) = \max_e c_e(f_e)\)
    \item \(\Phi(f) = \sum_{e} c_e(f_e) f_e\)
  \end{enumerate}
  \item The marginal cost of an edge with cost function \(c\) is ...
  \begin{enumerate}
    \item \(c^*(x) = \int_0^x c(t)\,dt\)
    \item \(c^*(x) = c(x)/x\)
    \item \(c^*(x) = \frac{d}{dx}\big(x\,c(x)\big)\)
    \item \(c^*(x) = c'(x)\)
  \end{enumerate}
  \item The cost of a path \(P\) at flow \(x\) is ...
  \begin{enumerate}
    \item the number of edges it contains
    \item the sum of the edge costs along it, \(c_P(x) = \sum_{e \in P} c_e(x)\)
    \item the cost of its cheapest edge
    \item the cost of its most expensive edge
  \end{enumerate}
  \item In Braess's paradox, adding a zero-cost road changes the Nash flow so that its cost ...
  \begin{enumerate}
    \item rises, even though a new option was added
    \item becomes zero
    \item always falls
    \item stays exactly the same
  \end{enumerate}
\end{enumerate}
\end{document}
